\begin{tabular}{lrlllll}
\toprule
model & seq\_len & a & b & c & d & e \\
\midrule
small & 256 & 前向 GPU 执行跨度为 21.62 ms，kernel 累计时间为 10.77 ms。该长度没有同配置的旧 Python 结果。 & 前向累计时间最多的是 ampere\_sgemm\_128x64\_tn，每次前向调用 85 次，占 kernel 时间 69.0\%。前向加反向的榜首为 ampere\_sgemm\_128x64\_tn，名称相同。 & 主要非矩阵乘法 kernel 包括 elementwise / MulFunctor / tensor × tensor（4.1\%）、vectorized / MulFunctor / tensor × tensor（3.0\%）、vectorized / CUDAFunctor\_add（2.5\%）；具体标签归属见 (c) 表。 & 完整训练的矩阵乘法占比为 54.7\%，前向为 74.4\%，变化 -19.7 个百分点。优化器阶段没有矩阵乘法，反向也增加了逐元素与归约等 kernel。 & 全部 MSA 层的 softmax 与矩阵乘法分别耗时 0.490、2.498 ms，时间比为 19.63\%。按本文标量运算约定，运算量比仅 0.28\%。 \\
small & 512 & 前向 GPU 执行跨度为 26.34 ms，kernel 累计时间为 24.05 ms。旧 Python 测得 25.85 ms，两者相差 +1.9\%。 & 前向累计时间最多的是 ampere\_sgemm\_128x64\_tn，每次前向调用 84 次，占 kernel 时间 55.3\%。前向加反向的榜首为 ampere\_sgemm\_128x64\_tn，名称相同。 & 主要非矩阵乘法 kernel 包括 vectorized / MulFunctor / scalar × tensor（5.2\%）、vectorized / MulFunctor / tensor × tensor（4.4\%）、vectorized / exp\_kernel\_cuda（3.9\%）；具体标签归属见 (c) 表。 & 完整训练的矩阵乘法占比为 51.0\%，前向为 65.4\%，变化 -14.4 个百分点。优化器阶段没有矩阵乘法，反向也增加了逐元素与归约等 kernel。 & 全部 MSA 层的 softmax 与矩阵乘法分别耗时 2.796、5.038 ms，时间比为 55.49\%。按本文标量运算约定，运算量比仅 0.49\%。 \\
small & 1024 & 前向 GPU 执行跨度为 80.72 ms，kernel 累计时间为 74.12 ms。该长度没有同配置的旧 Python 结果。 & 前向累计时间最多的是 ampere\_sgemm\_128x64\_tn，每次前向调用 84 次，占 kernel 时间 35.0\%。前向加反向的榜首为 void at::native::vectorized\_elementwise\_kernel<(int)4, at::native::BinaryFunctor<float, float, float, at::native::binary\_internal::MulFunctor<float>>, std::array<char *, (unsigned long)3>>(int, T2, T3)，与前向不同。 & 主要非矩阵乘法 kernel 包括 elementwise / masked\_fill\_kernel（7.0\%）、vectorized / MulFunctor / scalar × tensor（7.0\%）、vectorized / exp\_kernel\_cuda（7.0\%）；具体标签归属见 (c) 表。 & 完整训练的矩阵乘法占比为 39.0\%，前向为 46.2\%，变化 -7.2 个百分点。优化器阶段没有矩阵乘法，反向也增加了逐元素与归约等 kernel。 & 全部 MSA 层的 softmax 与矩阵乘法分别耗时 21.045、14.267 ms，时间比为 147.50\%。按本文标量运算约定，运算量比仅 0.78\%。 \\
medium & 256 & 前向 GPU 执行跨度为 42.82 ms，kernel 累计时间为 31.60 ms。该长度没有同配置的旧 Python 结果。 & 前向累计时间最多的是 ampere\_sgemm\_128x64\_tn，每次前向调用 169 次，占 kernel 时间 74.8\%。前向加反向的榜首为 ampere\_sgemm\_128x64\_tn，名称相同。 & 主要非矩阵乘法 kernel 包括 elementwise / MulFunctor / tensor × tensor（3.2\%）、vectorized / MulFunctor / tensor × tensor（3.1\%）、vectorized / CUDAFunctor\_add（2.1\%）；具体标签归属见 (c) 表。 & 完整训练的矩阵乘法占比为 52.7\%，前向为 78.1\%，变化 -25.4 个百分点。优化器阶段没有矩阵乘法，反向也增加了逐元素与归约等 kernel。 & 全部 MSA 层的 softmax 与矩阵乘法分别耗时 1.194、6.792 ms，时间比为 17.58\%。按本文标量运算约定，运算量比仅 0.22\%。 \\
medium & 512 & 前向 GPU 执行跨度为 81.81 ms，kernel 累计时间为 75.79 ms。旧 Python 测得 80.62 ms，两者相差 +1.5\%。 & 前向累计时间最多的是 ampere\_sgemm\_128x64\_tn，每次前向调用 168 次，占 kernel 时间 56.4\%。前向加反向的榜首为 ampere\_sgemm\_128x64\_tn，名称相同。 & 主要非矩阵乘法 kernel 包括 vectorized / MulFunctor / tensor × tensor（5.1\%）、vectorized / exp\_kernel\_cuda（4.6\%）、vectorized / MulFunctor / scalar × tensor（4.5\%）；具体标签归属见 (c) 表。 & 完整训练的矩阵乘法占比为 48.8\%，前向为 63.2\%，变化 -14.3 个百分点。优化器阶段没有矩阵乘法，反向也增加了逐元素与归约等 kernel。 & 全部 MSA 层的 softmax 与矩阵乘法分别耗时 10.992、15.548 ms，时间比为 70.70\%。按本文标量运算约定，运算量比仅 0.39\%。 \\
\bottomrule
\end{tabular}
